\subsection{LOCAL DEFINITION OF A LIE GROUP}

PROPOSITION 18. Let \(\mathbf{G}\) be a group and \(\mathbf{U}\) and \(\mathbf{V}\) two subsets of \(\mathbf{G}\) containing \(e\) . Suppose that \(\mathbf{U}\) has an analytic manifold structure satisfying the following conditions:

\begin{enumerate}
\def\labelenumi{(\roman{enumi})}
\item
  \(\mathrm{V} = \mathrm{V}^{-1},\mathrm{V}^2\subset \mathrm{U},\mathrm{V}\) is open in U;
\item
  the mapping \((x,y)\mapsto xy^{-1}\) of \(\mathbf{V}\times \mathbf{V}\) into \(\mathbf{U}\) is analytic;
\item
  for all \(g \in \mathbf{G}\) , there exists an open neighbourhood \(\mathbf{V}'\) of \(e\) in \(\mathbf{V}\) such that \(g\mathbf{V}'g^{-1} \subset \mathbf{U}\) and such that the mapping \(x \mapsto gxg^{-1}\) of \(\mathbf{V}'\) into \(\mathbf{U}\) is analytic.
\end{enumerate}

Then there exists one and only one analytic manifold structure on \(\mathbf{G}\) with the following properties:

(α) G with this structure is a Lie group;

(β) V is open in G;

\begin{enumerate}
\def\labelenumi{(\Alph{enumi})}
\setcounter{enumi}{24}
\tightlist
\item
  the manifold structures on G and U induce the same structure on V.
\end{enumerate}

\begin{enumerate}
\def\labelenumi{(\alph{enumi})}
\item
  Let A be an open subset of V and \(v_{0}\) an element of V such that \(v_{0}A \subset V\) . Then \(v_{0}A\) is the set of \(v \in V\) such that \(v_{0}^{-1}v \in A\) and hence is an open subset of V (taking account of (ii)). Moreover, (ii) implies that the mappings \(v \mapsto v_{0}v\) of A onto \(v_{0}A\) and \(v \mapsto v_{0}^{-1}v\) of \(v_{0}A\) onto A are inverse analytic bijections and hence analytic isomorphisms.
\item
  We choose an open neighbourhood \(W\) of \(e\) in \(V\) such that \(W = W^{-1}\) , \(W^3 \subset V\) and there exists a chart \((W, \phi, E)\) of the manifold \(U\) with domain \(W\) . For all \(g \in G\) , let \(\phi_g\) be the mapping \(h \mapsto \phi(g^{-1}h)\) of \(gW\) into \(E\) . We show that the charts \(\phi_g\) of \(G\) are analytically compatible. Let \(g_1, g_2\) be elements of \(G\) such that \(g_1W \cap g_2W \neq \emptyset\) , so that \(g_2^{-1}g_1\) and \(g_1^{-1}g_2\) belong to \(W^2\) . By (a), \(W \cap g_1^{-1}g_2W\) is an open subset of \(W\) and hence
\end{enumerate}

\[
\phi_ {g _ {1}} (g _ {1} \mathrm{W} \cap g _ {2} \mathrm{W}) = \phi (\mathrm{W} \cap g _ {1} ^ {- 1} g _ {2} \mathrm{W})
\]

is an open subset D of E. For \(d \in \mathbf{D}\) ,

\[
\left(\phi_ {g _ {2}} \circ \phi_ {g _ {1}} ^ {- 1}\right) (d) = \phi \left(g _ {2} ^ {- 1} g _ {1} \phi^ {- 1} (d)\right);
\]

by (a) we see that \(\phi_{g_2} \circ \phi_{g_1}^{-1}\) is analytic.

\begin{enumerate}
\def\labelenumi{(\alph{enumi})}
\setcounter{enumi}{2}
\item
  By (b) there exists on \(\mathbf{G}\) an analytic manifold structure such that \((\phi_g)_{g\in G}\) is an atlas on \(\mathbf{G}\) . For all \(g_0\in \mathbf{G}\) , the mapping \(g\mapsto g_0g\) ( \(g\in \mathbf{G}\) ) leaves this atlas invariant and hence is an automorphism of \(\mathbf{G}\) with this manifold structure. In particular condition \((\mathrm{GL}_1)\) is satisfied.
\item
  Let \(v_{0} \in V\) . By (ii) there exists an open neighbourhood \(A\) of \(e\) in \(W\) such that \(v_{0}A \subset V\) . This proves first that \(V\) is open in \(G\) . By (a) the mapping \(v \mapsto v_{0}v\) of \(A\) onto \(v_{0}A\) is an analytic isomorphism with the structures induced by U. By (c) we see that the manifold structures on G and U induce the same structure on \(v_{0}\mathbf{A}\) and hence finally on V.
\item
  by (d), (ii) and (iii), we see that conditions \((\mathrm{GL}_2)\) and \((\mathrm{GL}_3)\) hold. Hence G is a Lie group.
\item
  If a manifold structure on G is compatible with the group structure of G and is such that V is an open submanifold of G, then \((\phi_{g})_{g\in G}\) is an atlas on G. Hence the uniqueness assertion of the proposition.
\end{enumerate}

PROPOSITION 19. Let G be a topological group, H a Lie group and f a homomorphism of the group G into the group H. Suppose that there exist an open neighbourhood of \(e_{G}\) in G, a chart (V, \(\phi\) , E) of the manifold H at \(e_{H}\) and a closed vector subspace F of E admitting a topological supplement, such that \(f(\mathrm{U}) \subset \mathrm{V}\) and \((\phi \circ f)\mid_U\) is a homeomorphism of U onto \(\phi(\mathrm{V}) \cap \mathrm{F}\) . Then there exists a unique manifold structure on G such that f is an immersion; this structure is the inverse image under f of the manifold structure on H. With this structure G is a Lie group.

As translations of G (resp. H) are homeomorphisms (resp. analytic isomorphisms), f satisfies condition (R) of Differentiable and Analytic Manifolds, R, 5.8.1. The first two assertions of the proposition then follow from Differentiable and Analytic Manifolds, R, loc. cit. Consider the commutative diagram

\[
\begin{array}{c} \mathrm{G} \times \mathrm{G} \xrightarrow {m} \mathrm{G} \\ f \times f \Big \downarrow \qquad \qquad \qquad \qquad \qquad \qquad \qquad \qquad \qquad \Big \downarrow f \\ \mathrm{H} \times \mathrm{H} \xrightarrow {n} \mathrm{H} \end{array}
\]

where \(m(x,y)=xy^{-1}\) (resp. \(n(x,y)=xy^{-1}\) ) for x, y in G (resp. H). Then \(n\circ(f\times f)\) is analytic, hence \(f\circ m\) is analytic and hence m is analytic since f is an immersion. Therefore G is a Lie group.

The Lie group structure on G is called the inverse image of the Lie group structure on H under f.

COROLLARY. Let G be a topological group, N a discrete normal subgroup of G and \(\pi\) the canonical mapping of G onto G/N. Suppose that an analytic manifold structure is given on G/N, compatible with the topological group structure on G/N. Then there exists a unique manifold structure on G such that \(\pi\) is an immersion; this structure is the inverse image under \(\pi\) of the manifold structure on G/N. With this structure, \(\pi\) is étale, G is a Lie group and G/N is the quotient Lie group of G by N.

Remark. Let H be a connected real or complex Lie group, \(\hat{H}\) its universal covering \(\dagger\) and \(\pi\) the canonical mapping of \(\hat{H}\) onto H. When we speak of \(\hat{H}\) as a Lie group, we shall always mean with the inverse image structure of that on \(\mathbf{H}\) under \(\pi\) .

